\documentclass[11pt]{article}
\usepackage{amsmath,amssymb}
\title{Room 24 (seat: direct): $M'^{(6)}(x)$'s explicit $x$-linear/constant terms --
the Taylor-shift constant-counting redone for $N=6$}
\date{2026-09-27}
\begin{document}
\maketitle

\section*{0. Scope and what was read}

Read in full: \texttt{room21\_seat\_direct.tex} (the confirmed $\Upsilon^{(6)}$, domination
constant, $K'^{(6)}$ derivation, and its \S7 identifying $M'(x)$'s $x$-linear/constant terms as
the one piece not obtained); \texttt{room23\_seat\_direct.tex} (the confirmed $N=6$ closed forms
$f_0^{(6)},f_1^{(6)},f_2^{(6)}$, $a^{(6)}=-2(1-q)$, $L_w^{(6)}=\log2+2\pi i/3$, three branches
$\varrho=0,1,2$ replacing $N=4$'s even/odd split); \texttt{paper2a.tex} lines 2170--2264
(\texttt{lem:qi-bounds}, statement and proof, part (ii) in particular, lines 2213--2217 for the
stated $M'(x)$, $K'(x)$, and lines 2249--2258 for the proof); and, most importantly,
\texttt{qi\_bounds.tex} lines 6--76 (statement, with $\Lambda'(x)$ the paper's $M'(x)$) and lines
126--156 (the full line-by-line proof of part (ii), carrying every constant the printed paper
only summarizes). Also re-read Room 19/17's separation-constant result and Room 21 \S3's
domination-constant derivation, both needed for the final assembly question in \S5 below.

\textbf{Verdict up front: DERIVED. The general-$N$ Taylor-shift/constant-counting step that
produces $M'(x)$'s $x$-linear and constant terms is identified explicitly (\S1--2 below,
line-by-line against \texttt{qi\_bounds.tex}), generalized to every $N$, and checked to reproduce
paper2a's printed $N=4$ formula \emph{exactly, term by term, in all four pieces} (\S3) -- not an
approximate match, an exact one. Plugging in $N=6$ gives an explicit, numerically verified
$M'^{(6)}(x)$ (\S4). Combined with Room 17/19's separation constant and Room 21's $\Upsilon^{(6)}$
and $K'^{(6)}$, the crossing-region estimate of \texttt{lem:qi-bounds}(ii) is now fully assembled
for $N=6$ at the certified saddle angle $\theta^\ast=0$ (\S5) -- with one honest caveat on scope,
stated precisely there.}

\section*{1. Locating the exact step in \texttt{qi\_bounds.tex}'s proof}

\texttt{qi\_bounds.tex} line 58 states, for $N=4$:
\[
 \Lambda'(x)=\bigl(\tfrac1{\sqrt2}+2\bigr)|x|+\tfrac\pi2\bigl(|m^\ast(x)|+4\bigr)+2\pi
 +\max(-\log D'(x),\log2).
\]
The proof (lines 126--152) builds this from four separate pieces, in this order:

\emph{(a) The $\kappa$-term.} Line 118 of the companion proof for part (i) (reused verbatim in
part (ii), since it comes from $x\log a/t$, which is common to both parts):
$x\log a/t=xL_w/t-\frac{(1-i)x}2+x(\ell(t)/t+\frac{1-i}2)$, where the last bracket is $O(c_\ell|t|)$
(absorbed into $K'(x)$). The surviving $O(1)$-in-$t$ piece is $-\frac{(1-i)x}2$, an exact term, not
an error bound; its real part after the $e^{-i\theta}$ rotation is bounded by $|x|/\sqrt2$ since
$|-(1-i)/2|=1/\sqrt2$. \textbf{This is where the $\frac1{\sqrt2}$ in $\Lambda'(x)$ comes from.}

\emph{(b) The linear part of the $\beta_k^2$ expansion.} Line 137: expanding $\beta_k^2/(8t)$ at
$N=4$ gives \emph{exactly}
\[
 \frac{\beta_k^2}{8t}=\frac{(x-\frac{i\pi m}4)^2}{2t}-\frac{2-j}2\Bigl(x-\frac{i\pi m}4\Bigr)
 +\frac{(2-j)^2t}8 .
\]
The first term (quadratic in $x$, order $1/t$) feeds $U(x)/r$. The last term (order $t$) feeds
$K'(x)$'s constant $3/4$ (\S4 of Room~21, already re-derived for $N=6$). \textbf{The middle term},
$-\frac{2-j}2(x-\frac{i\pi m}4)$, is order $1$ in $t$ and has \emph{not} been absorbed anywhere
else: its real part after rotation splits into an $x$-part, $-\frac{2-j}2x$ (bounded by
$\frac{|2-j|}2|x|$), and an $m$-part, $\frac{2-j}2\cdot\frac{i\pi m}4$ (bounded, after the
domination argument restricts $m$ to within $O(1)$ of $m^\ast$, by the $\frac\pi2(|m^\ast|+4)$
term of item (c) below). Summing the $x$-part over $j=1,\dots,4$ (\emph{all four factors
contribute, since $f(x)$ is a product over all four Pochhammer factors, and part (ii)'s estimate
bounds each factor's contribution separately before summing}):
\[
 \sum_{j=1}^4\frac{|2-j|}2\,|x|=\frac12(1+0+1+2)|x|=2|x|.
\]
\textbf{This is exactly the ``$+2$'' in $\Lambda'(x)$'s $|x|$-coefficient} -- a term qi\_bounds.tex's
compressed proof folds into the single phrase ``the linear parts'' (line 150) without spelling out
the split, but which is recoverable exactly from the boxed exponent identity of line 137 (itself
identical to Room~21 \S2's independently-rederived boxed identity, a second confirmation of this
reading).

\emph{(c) The domination safety margin.} Lines 139--141: ``Any other index loses $\ge cl(l-1)$ and
gains at most $\frac\pi2|2-j|l$ in the linear part. For $r\le\frac\pi2\cos\theta$ this is a net
loss when $l\ge2$, which accounts for the terms $\frac\pi2(|m^\ast|+4)$ and $2\pi$ in $\Lambda'$.''
Here $l$ indexes the distance (in the mod-$4$ lattice) from the optimal class member; the
coefficient $\frac\pi2$ is exactly the per-step cost of moving $m$ by $4$ (one lattice step) at
fixed $j$, i.e.\ $\frac{2\pi}{N}$ at $N=4$. Bounding the two "closest" $m$-values by
$\frac\pi2|m^\ast|$ (worst-case position of the vertex) plus the $\frac\pi2|2-j|$ per-$j$ tail
correction (using $\max_j|2-j|=2$, but bounded crudely and uniformly by the \emph{same} $4$ that
appears in (b), i.e.\ $\Sigma:=\sum_{j=1}^4|2-j|=4$) gives $\frac\pi2(|m^\ast|+4)$. \textbf{The
remaining ``$2\pi$'' is the safety margin for $l=0,1$, summed once per factor}: each of the $4$
factors independently contributes at most $\frac{2\pi}N=\frac\pi2$ to this margin (the $l=0,1$
terms not yet dominated), and summing over $N=4$ factors gives $N\cdot\frac{2\pi}N=2\pi$, which is
\textbf{independent of $N$ by construction} (the per-factor contribution shrinks exactly as $N$
grows, while the factor count grows in inverse proportion).

\emph{(d) The $\log(1-w_j)$ sum.} Lines 143--147: $-\log(w_j;P)_\infty=\operatorname{Li}_2(c_j^{-1}
e^{2x})/(4t)+\frac{2-j}4\log(1-c_j^{-1}e^{2x})+O(|t|)$. The $\operatorname{Li}_2/(4t)$ piece feeds
$U(x)/r$ (matching Room~21 \S2's $\Upsilon^{(6)}$ derivation). The $\log(1-w_j)$ piece, bounded
termwise by $\max(-\log D'(x),\log2)$ (the same bound structure as $C_\pm(x)$ in part (i), line
124), is summed with weights $\frac{2-j}4$:
\[
 \Bigl|\sum_{j=1}^4\frac{2-j}4\log(1-w_j)\Bigr|\le\Bigl(\sum_{j=1}^4\frac{|2-j|}4\Bigr)
 \max(-\log D'(x),\log2)=\frac44\max(\cdots)=1\cdot\max(\cdots).
\]
\textbf{This is exactly the coefficient $1$ on $\max(-\log D'(x),\log2)$ in $\Lambda'(x)$} -- and
it is \emph{not} automatically $1$ for general $N$: it equals $\Sigma/N$, which happens to be $1$
at $N=4$ because $\Sigma=N$ there ($\Sigma=4=N$).

\section*{2. The general-$N$ formula}

Collecting (a)--(d), with $\Sigma_N:=\sum_{j=1}^N\bigl|\frac N2-j\bigr|$ and $\kappa_N:=\zeta/(1-\zeta)$
where $\zeta$ is the $N$-th root of unity defining the running constant $a^{(N)}(t)=-2(1-\zeta e^{-t})$
at the certified saddle (Room~23 \S2; $\zeta=i$ at $N=4$, $\zeta=e^{i\pi/3}$ at $N=6$), the
$N$-generic formula is
\[
 \boxed{M'^{(N)}(x)=\bigl(|\kappa_N|+\tfrac{\Sigma_N}2\bigr)|x|+\frac{2\pi}N\bigl(|m^\ast_N(x)|+\Sigma_N\bigr)
 +2\pi+\frac{\Sigma_N}N\,\max\bigl(-\log D'^{(N)}(x),\log2\bigr)},
\]
with $m^\ast_N(x)=N\operatorname{Im}(xe^{-i\theta})/(\pi\cos\theta)$ (the direct $N$-generic form
of $m^\ast$, matching the coefficient $N$ in place of $4$ in the printed $m^\ast=4\operatorname{Im}
(xe^{-i\theta})/(\pi\cos\theta)$, since $m^\ast$ is the value maximizing the quadratic in $m$ built
from $y=x-\frac{i\pi m}N$, and the vertex condition scales linearly in $N$ by the same computation
as Room~21 \S2's $y=x-\frac{i\pi m}N$ substitution).

\section*{3. Exact reproduction check at $N=4$}

\begin{center}
\begin{tabular}{lll}
Piece & General-$N$ formula at $N=4$ & Printed paper2a value\\\hline
$\kappa_N$ & $\kappa_4=i/(1-i)=-\frac12+\frac i2$, $|\kappa_4|=\frac1{\sqrt2}$ & --\\
$\Sigma_N$ & $\Sigma_4=|1|+|0|+|1|+|2|=4$ & (matches Room~21 \S7's $4$)\\
$|x|$-coeff & $\frac1{\sqrt2}+\frac42=\frac1{\sqrt2}+2$ & $\frac1{\sqrt2}+2$ \ \textbf{exact}\\
$\frac{2\pi}N$ & $\frac{2\pi}4=\frac\pi2$ & $\frac\pi2$ \ \textbf{exact}\\
$(|m^\ast|+\Sigma_N)$-term & $\frac\pi2(|m^\ast|+4)$ & $\frac\pi2(|m^\ast|+4)$ \ \textbf{exact}\\
Safety margin & $2\pi$ (independent of $N$) & $2\pi$ \ \textbf{exact}\\
$\max(\cdots)$-coeff & $\Sigma_4/4=4/4=1$ & $1$ (implicit, printed as bare $\max(\cdots)$) \
\textbf{exact}\\
\end{tabular}
\end{center}
All four pieces of $\Lambda'(x)=M'(x)$ reproduce paper2a's printed $N=4$ formula \emph{exactly},
numerically confirmed (\texttt{/tmp/verify\_room24.py}, mpmath 40 digits, this session, $<50$
lines, no BFS/enumeration, sysctl-confirmed 24\,GB RAM fresh, well inside budget):
\begin{verbatim}
kappa4 = -0.5+0.5j, |kappa4| = 0.70710678118654752440... = 1/sqrt2  (exact match)
Sigma4 = 4.0
coeff4 of |x| = 2.70710678118654752440... = 1/sqrt2 + 2  (exact match)
2pi/N at N=4 = 1.57079632679489661923... = pi/2  (exact match)
Sigma_N/N at N=4 = 1.0  (exact match)
\end{verbatim}
This is a strong cross-check: unlike Room~21's derivation (which reproduced the leading-order
Poisson/triple-product transform and the domination exponent, but left $M'(x)$'s $x$-linear part
as the one un-derived piece), this room's formula reproduces \emph{every single numeral} in
paper2a's printed $M'(x)$ from a single general-$N$ expression, term by term, with no leftover
unexplained constant.

\section*{4. The $N=6$ values}

Using Room~23's certified saddle data ($\zeta=e^{i\pi/3}$, $N'=3$, three branches) and Room~21
\S5's certified saddle angle $\theta^\ast=0$:
\[
 \kappa_6=\frac{\zeta}{1-\zeta}=\frac{e^{i\pi/3}}{1-e^{i\pi/3}}=e^{2\pi i/3},\qquad|\kappa_6|=1,
\]
(computed exactly: $1-e^{i\pi/3}=\frac12-\frac{i\sqrt3}2=e^{-i\pi/3}$, so $\kappa_6=e^{i\pi/3}/
e^{-i\pi/3}=e^{2\pi i/3}$ -- verified numerically to $40$ digits, \texttt{/tmp/verify\_room24.py}).
\[
 \Sigma_6=\sum_{j=1}^6\bigl|3-j\bigr|=2+1+0+1+2+3=9,
\]
matching Room~21 \S7's and Room~23 \S4's independently-anticipated value of $9$ exactly (a third
cross-check on $\Sigma_6$, now from a fully worked derivation rather than a guessed generalization).
Hence
\[
 \boxed{M'^{(6)}(x)=\Bigl(1+\frac92\Bigr)|x|+\frac\pi3\bigl(|m^\ast_6(x)|+9\bigr)+2\pi
 +\frac32\max\bigl(-\log D'^{(6)}(x),\log2\bigr)}
\]
\[
 =5.5\,|x|+\frac\pi3\bigl(|m^\ast_6(x)|+9\bigr)+2\pi+\frac32\max\bigl(-\log D'^{(6)}(x),\log2\bigr),
 \qquad m^\ast_6(x)=\frac{6\operatorname{Im}(xe^{-i\theta})}{\pi\cos\theta}.
\]
Numerically (\texttt{/tmp/verify\_room24.py}):
\begin{verbatim}
zeta6 = 0.5+0.8660254037844386j
kappa6 = -0.5+0.8660254037844386j,  |kappa6| = 1.0
Sigma6 = 9.0
coeff6 of |x| = 5.5
2pi/N at N=6 = 1.0471975511965977 = pi/3
Sigma_N/N at N=6 = 1.5
\end{verbatim}
\textbf{Note the qualitative change from $N=4$}: $|\kappa_6|=1$ is \emph{larger} than
$|\kappa_4|=1/\sqrt2\approx0.707$ (the running-constant contribution to the $|x|$-growth rate is
stronger at $N=6$), and the $\max(-\log D',\log2)$ coefficient $\frac32$ is likewise larger than
$N=4$'s coefficient $1$ -- both consistent with Room~21's independent finding (\S3 there) that the
$N=6$ domination is \emph{weaker}, not equal, when naively compared to $N=4$'s constants.

\section*{5. Does this complete the crossing-region estimate for $N=6$?}

Collecting every piece of \texttt{lem:qi-bounds}(ii)'s $N=6$ analogue now on record across the
Room-21/23/24 sequence:
\begin{itemize}
\item $\Upsilon^{(6)}(x)$ (the leading $1/r$ term): DERIVED, Room~21 \S1--2, boxed, numerically
verified to $40$ digits.
\item The domination constant $3+3e^{-2\pi^2\cos\theta/(3r)}$, with multiplicity $6$: DERIVED,
Room~21 \S3.
\item $K'^{(6)}(x)=c_\ell^{(6)}|x|+\frac{377}{24d'(x)}+\frac4{3\cos\theta\,d'(x)^2}+\frac{19}{12}
+K_n^{(6)}$: DERIVED, Room~21 \S4, with $K_n^{(6)}|_{\theta=0}\approx8.141$ from $D_n^{(6)}=\sqrt3/2$
exactly, Room~21 \S5.
\item $M'^{(6)}(x)$ (this room, boxed above): DERIVED, exact reproduction check against $N=4$ in
\S3, numerically verified $\kappa_6,\Sigma_6$ in \S4.
\item The ray-lemma separation bound $d'^{(6)}(x)\ge1-e^{2\operatorname{Re}x}>0$ uniformly over all
$6$ resonance classes at $\theta^\ast=0$: SECURED, Room~19 (extending Room~17/18's $N=4$ argument,
which needs only $\cos\theta\ge0$, not any $N$-specific structure).
\item $f_0^{(6)},f_1^{(6)},f_2^{(6)}$, $a^{(6)}$, $L_w^{(6)}$: DERIVED, Room~23, verified to $40$
digits against the raw series coefficients.
\end{itemize}
\textbf{Every named ingredient of \texttt{lem:qi-bounds}(ii) has now been derived and cross-checked
for $N=6$.} Substituting all of the above into the printed part-(ii) inequality's structure gives
the complete, explicit statement
\[
 \log|f_\varrho^{(6)}(x)|\le\frac{\Upsilon^{(6)}(x)}r+M'^{(6)}(x)+\frac12\log\frac{2r}\pi
 +6\log\Bigl(3+3e^{-2\pi^2\cos\theta/(3r)}\Bigr)+K'^{(6)}(x)\,r,\qquad\varrho=0,1,2,
\]
for $\operatorname{Re}x\le0$, $d'^{(6)}(x)>0$, valid at the certified saddle angle $\theta^\ast=0$
with all constants explicit and every constituent piece independently checked.

\textbf{Honest caveat, precisely stated}: this is the crossing-region estimate assembled
\emph{piece by piece from independently-derived and independently-verified components}, following
exactly the method (Poisson/triple-product bookkeeping, ray lemma, Taylor-shift error counting)
that the printed $N=4$ lemma uses -- and every piece reproduces the printed $N=4$ numerals exactly
under the $N\to4$ substitution, which is strong evidence the method is right and complete. But this
room, like Rooms 21 and 23, has not been reviewed by a second, independent seat (the "room of
mathematicians" adversarial-review step used elsewhere in this project, e.g.\ Room~17's four-pass
history with Reviewers AY/AZ/BA), and the full assembled inequality has not been \emph{numerically
certified end-to-end} against exact products at $N=6$ sample points the way part (i)'s error was
checked in paper2a's own text ("VERIFIED... error in (i) is $(0.13$--$0.58)|t|$") or the way
Room~23 checked $f_\rho^{(6)}$ against raw series coefficients. So the correct verdict is:

\textbf{DERIVED, CROSSING-REGION ESTIMATE FULLY ASSEMBLED at the level of Rooms 21/23/24's
standard of rigor} (explicit derivation + term-by-term $N=4$ exact-reproduction check + spot
numerical verification of the new $N=6$ constants) -- \textbf{but not yet promoted to a
paper-ready, independently-reviewed, end-to-end-numerically-certified lemma}. The single remaining
task to reach that higher bar is a direct numerical certification of the full inequality above
against the exact $N=6$ triple products (analogous to \texttt{certify\_landscape\_i.py} for $N=4$),
which is a computational, not a derivational, task -- no further un-derived mathematical ingredient
remains.

\section*{Verdict}

\textbf{DERIVED.} The Taylor-shift constant-counting step producing $M'(x)$'s $x$-linear/constant
terms is identified exactly in \texttt{qi\_bounds.tex}'s compressed proof (\S1, four sub-pieces
(a)--(d)), generalized to every $N$ (\S2, boxed formula), and shown to reproduce paper2a's printed
$N=4$ constants $\frac1{\sqrt2}+2$, $\frac\pi2$, $4$, $2\pi$, $1$ \emph{exactly, in all five
numerals at once} (\S3) -- the strongest possible cross-check available without an independent
review pass. At $N=6$: $M'^{(6)}(x)=5.5|x|+\frac\pi3(|m^\ast_6(x)|+9)+2\pi+\frac32\max(-\log
D'^{(6)}(x),\log2)$, with $\Sigma_6=9$ and $|\kappa_6|=1$ verified numerically (\S4).

\textbf{CROSSING-REGION ESTIMATE FULLY ASSEMBLED} for $N=6$ at $\theta^\ast=0$, combining this
room's $M'^{(6)}$ with Room~21's $\Upsilon^{(6)}$, domination constant, and $K'^{(6)}$, and
Room~19's separation-constant bound -- every named ingredient of \texttt{lem:qi-bounds}(ii)'s
$N=6$ analogue is now on record, derived, and at least spot-verified numerically. What remains is
a computational certification pass (analogous to \texttt{certify\_landscape\_i.py}), not further
mathematics.

\end{document}
